\documentclass[../../../include/open-logic-section]{subfiles}

\begin{document}

\olfileid[pt]{sth}{ord-arithmetic}{mult}

\olsection{Multiplicação Ordinal}

Passamos agora à multiplicação ordinal, abordando-a de forma muito
semelhante à adição ordinal. Suponha-se que se pretende multiplicar $\alpha$
por~$\beta$. Pode imaginar-se uma grelha retangular, com largura $\alpha$
e altura $\beta$; o produto de $\alpha$ e $\beta$ é então o
resultado de percorrer cada linha e depois passar à linha seguinte\ldots\
até percorrer toda a grelha. Dito de outro modo,
o produto de $\alpha$ e $\beta$ obtém-se substituindo \emph{cada}
elemento de $\beta$ por uma cópia de $\alpha$.

Para formalizar, usa-se simplesmente a ordem lexicográfica inversa
no produto cartesiano de $\alpha$ e $\beta$:

\begin{defn}
Para quaisquer ordinais $\alpha, \beta$, o seu produto $\alpha \ordtimes \beta = \ordtype{\alpha \times \beta, \rlexless}$.
\end{defn}
\noindent
É preciso de novo confirmar que esta definição está bem formada:

\begin{lem}\ollabel{ordtimeslessiswo}
$\tuple{\alpha \times \beta, \rlexless}$ é um bom ordenamento, para quaisquer
ordinais $\alpha$ e $\beta$.
\end{lem}

\begin{proof}
Exatamente como para \olref[add]{ordsumlessiswo}.
\end{proof}
\noindent
E não é difícil provar que a multiplicação se comporta assim:

\begin{lem}\ollabel{ordtimesrecursion}
Para quaisquer ordinais $\alpha, \beta$:
\begin{align*}
	\alpha \ordtimes 0 &= 0\\
	\alpha \ordtimes (\beta \ordplus 1) &=
		(\alpha \ordtimes \beta) \ordplus \alpha\\
	\alpha  \ordtimes \beta &=
		\supstrict_{\delta < \beta}(\alpha \ordtimes \delta) &&
		\text{quando $\beta$ é um ordinal limite}.
\end{align*}
\end{lem}

\begin{proof}
Deixado como exercício.
\end{proof}

De fato, tal como no caso da adição, poderíamos ter definido a
multiplicação ordinal por estas equações recursivas, em vez de dar uma
definição direta. Tal como na adição, certo comportamento é
familiar:

\begin{lem}\ollabel{ordinalmultiplicationisnice}
Se $\alpha, \beta, \gamma$ são ordinais, então:
\begin{enumerate}
	\item\ollabel{ordtimes1} se $\alpha \neq 0$ e $\beta < \gamma$,
	então $\alpha \ordtimes \beta < \alpha \ordtimes \gamma$;
	\item\ollabel{ordtimes2} se $\alpha \neq 0$ e $\alpha \ordtimes
	\beta = \alpha\ordtimes\gamma$, então $\beta = \gamma$;
	\item\ollabel{ordtimes3}  $\alpha \ordtimes (\beta \ordtimes
	\gamma) = (\alpha \ordtimes \beta) \ordtimes \gamma$;
	\item\ollabel{ordtimes4}  Se $\alpha \leq \beta$, então $\alpha
	\ordtimes \gamma \leq \beta \ordtimes \gamma$;
	\item\ollabel{ordtimes5}  $\alpha \ordtimes (\beta \ordplus
	\gamma) = (\alpha \ordtimes \beta )\ordplus (\alpha\ordtimes
	\gamma)$.
\end{enumerate}
\end{lem}

\begin{proof}
Deixado como exercício.
\end{proof}

Pode provar-se (ou consultar) outros resultados, à vontade.
Mas, dado
\olref[ord-arithmetic][add]{ordsumnotcommute}, o
que se segue não deve surpreender:

\begin{prop}
A multiplicação ordinal não é comutativa: $2 \ordtimes \omega  =
\omega < \omega \ordtimes 2$
\end{prop}

\begin{proof}
$2 \ordtimes \omega = \supstrict_{n < \omega}(2\ordtimes  n) = \omega \in \supstrict_{n < \omega}(\omega \ordplus n) = \omega \ordplus \omega = \omega \ordtimes 2$.
\end{proof}
\noindent
De novo, a justificação intuitiva é direta. Para calcular $2
\ordtimes \omega$, substitui-se cada número natural por duas entidades.
Obtém-se o mesmo tipo de ordem se simplesmente se inserirem todos os
números ``meios'' nos naturais, isto é, se se considerar a
ordem natural em $\Setabs{\nicefrac{n}{2}}{n \in \omega}$. Assim posto, o tipo de ordem é claramente o mesmo que o de $\omega$
em si. Mas, para calcular $\omega \ordtimes 2$, dispõem-se duas
cópias de $\omega$, uma a seguir à outra.

\begin{prob}
Prove
\olref[sth][ord-arithmetic][mult]{ordtimeslessiswo},
\olref[sth][ord-arithmetic][mult]{ordtimesrecursion},
e
\olref[sth][ord-arithmetic][mult]{ordinalmultiplicationisnice}
\end{prob}

\end{document}
